\documentclass{beamer}
\usepackage{amsmath, amssymb, amsthm}
\usepackage{mathtools}
\usepackage{physics}
\usepackage{braket}
\usepackage{algorithm}
\usepackage{algpseudocode}
\usepackage{booktabs}
\usepackage{hyperref}
\usepackage{graphicx}
\usepackage{multirow}
\usepackage{array}

% Beamer theme
\usetheme{Madrid}
\usecolortheme{seahorse}
\setbeamertemplate{navigation symbols}{}

% Theorem environments for beamer
\setbeamertemplate{theorem}[ams style]
\setbeamertemplate{theorems}[numbered]

\theoremstyle{definition}
\newtheorem{defi}{Definition}[section]

\theoremstyle{plain}
\newtheorem{thm}{Theorem}[section]
\newtheorem{lem}[thm]{Lemma}
\newtheorem{cor}[thm]{Corollary}
\newtheorem{prop}[thm]{Proposition}

% Title page
\title[Week 2 Report]{Week 2 Internship Report: \\ Quantum Communication Theory \\ \small Noisy Quantum Channels and Information-Theoretic Limits}
\subtitle{Connection to Week 1: Foundations and Framework}
\author{Kaumud Sharma}
\institute{}
\date{}

\begin{document}

\begin{frame}
    \titlepage
\end{frame}

\begin{frame}{Outline}
    \tableofcontents
\end{frame}

\section{Introduction: Connection to Week 1 Foundations}

\begin{frame}{Week 1: Established Foundation}
    \begin{block}{Week 1 Achievements}
        \begin{enumerate}
            \item \textbf{Electromagnetic foundations}: Maxwell's equations, modal decomposition
            \item \textbf{Quantum field theory}: Field quantization, coherent states
            \item \textbf{Atom-field interaction}: Hamiltonian $H(t|\theta)$ with fiber parameters
            \item \textbf{Cq channel theory}: Encoding, Holevo capacity
            \item \textbf{Computational complexity}: NP-hardness via SJIP reduction
        \end{enumerate}
    \end{block}
    
    \begin{block}{Week 2: Extensions to Noisy Channels}
        \begin{enumerate}
            \item \textbf{Noise modeling}: GKSL master equation
            \item \textbf{Exact dynamics}: Solution methods
            \item \textbf{Noisy capacity}: Holevo capacity under noise
            \item \textbf{Fundamental limits}: Quantum DPI
            \item \textbf{Decoding ambiguity}: Formal characterization
        \end{enumerate}
    \end{block}
\end{frame}

\section{Mathematical Modeling of Fiber Noise}

\begin{frame}{From Ideal to Noisy Propagation}
\footnotesize

\begin{block}{Week 1: Ideal Model}
\textbf{Coherent State:} $|\alpha(0)\rangle \to |\alpha(L)\rangle$ with $\alpha(L) = \alpha(0)e^{i\beta L}e^{-\alpha L/2}$
\begin{itemize}
    \item Noiseless, perfect information, deterministic
\end{itemize}
\end{block}

\vspace{0.2cm}

\begin{block}{Week 2: Realistic Model with GKSL Master Equation}
\[
\dv{\rho}{t} = -i[H,\rho] + \sum_k \gamma_k\left(L_k \rho L_k^\dagger - \frac{1}{2}\{L_k^\dagger L_k, \rho\}\right)
\]

\textbf{Fiber Noise Sources:}
\begin{itemize}
    \item Amplitude damping: $\gamma_{\text{loss}}$, Phase damping: $\gamma_\phi$
    \item Non-unitary evolution, information loss, stochastic processes
\end{itemize}
\end{block}

\vspace{0.2cm}

\begin{alertblock}{Key Connection}
Week 1's attenuation $\alpha$ becomes Week 2's damping: $\gamma_{\text{loss}} = \frac{\alpha v_g \ln(10)}{10}$

Ideal model = special case when $\gamma_{\text{loss}} = \gamma_\phi = 0$
\end{alertblock}

\end{frame}

\begin{frame}{GKSL Master Equation}
    \begin{thm}[GKSL Master Equation]
        \[
        \dv{\rho_S}{t} = -i[H_S, \rho_S] + \sum_{k} \gamma_k \left( L_k \rho_S L_k^\dagger - \frac{1}{2} \{L_k^\dagger L_k, \rho_S\} \right)
        \]
    \end{thm}
    
    \begin{block}{Connection to Week 1 Parameters}
        \begin{align*}
            \gamma_{\text{loss}} &= \frac{\alpha v_g \ln(10)}{10} \\
            \gamma_\phi &= \frac{(\omega D_p \Delta L)^2 v_g}{2}
        \end{align*}
        where $\alpha$, $D_p$ are Week 1 fiber parameters
    \end{block}
\end{frame}

\section{Exact Solution Methods}

\begin{frame}{Direct Diagonalization Method}
    \begin{thm}[Direct Diagonalization]
        For $d$-dimensional system, Liouvillian $\mathbb{L}$ is $d^2 \times d^2$:
        \[
        \vec{\rho}(t) = e^{\mathbb{L}t} \vec{\rho}(0)
        \]
    \end{thm}
    
    \begin{block}{Complexity}
        \begin{itemize}
            \item Scales as $O(d^6)$
            \item Limited to small systems ($d \lesssim 10$)
            \item Exact up to numerical precision
        \end{itemize}
    \end{block}
\end{frame}

\begin{frame}{Quantum Trajectory Method}
    \begin{block}{For Non-Markovian Dynamics}
        \[
        \dv{}{t} \ket{\psi(t)} = -i H_{\text{eff}}(t) \ket{\psi(t)} + \sum_k \xi_k(t) L_k \ket{\psi(t - \tau_k)}
        \]
    \end{block}
    
    \begin{block}{Advantages}
        \begin{itemize}
            \item Handles large systems
            \item Captures exact noise correlations
            \item Parallelizable via Monte Carlo
        \end{itemize}
    \end{block}
\end{frame}

\section{Holevo Capacity Calculation}

\begin{frame}{From Ideal to Noisy Capacity}
    \begin{columns}[T]
        \begin{column}{0.48\textwidth}
            \begin{block}{Week 1: Ideal Capacity}
                \centering
                \[
                C(\theta) = \max_{p(x)} \chi(\{p(x), \rho(x|\theta)\})
                \]
                For ideal channels
            \end{block}
        \end{column}
        
        \begin{column}{0.48\textwidth}
            \begin{block}{Week 2: Noisy Capacity}
                \centering
                \[
                C_\chi(\mathcal{N}) = \max_{\{p(x), \rho_x\}} \chi(\{p(x), \mathcal{N}(\rho_x)\})
                \]
                For noisy channel $\mathcal{N}$
            \end{block}
        \end{column}
    \end{columns}
\end{frame}

\begin{frame}{Analytical Capacities for Specific Channels}
    \begin{thm}[Depolarizing Channel]
        \[
        C_\chi^{\text{dep}} = 1 - H_2\left( \frac{p}{2} \right) - p \log_2 3
        \]
    \end{thm}
    
    \begin{thm}[Fiber Noise Channel]
        \[
        C_\chi^{\text{fiber}} \leq \min \left\{ C_\chi^{\text{AD}}(\gamma), C_\chi^{\text{PD}}(\lambda) \right\}
        \]
    \end{thm}
    
    \begin{block}{For typical fiber parameters}
        $C_\chi \approx 0.7-0.9$ bits per photon
    \end{block}
\end{frame}

\section{Post-Channel CPTP Correction}

\begin{frame}{Quantum Data Processing Inequality}
    \begin{thm}[Quantum DPI]
        For $\rho \to \sigma \to \tau$ with $\tau = \mathcal{C}(\sigma)$:
        \[
        I(\rho; \tau) \leq I(\rho; \sigma)
        \]
    \end{thm}
    
    \begin{thm}[Capacity Non-Increase]
        For any $\mathcal{N}$ and CPTP $\mathcal{C}$:
        \[
        C_\chi(\mathcal{C} \circ \mathcal{N}) \leq C_\chi(\mathcal{N})
        \]
    \end{thm}
\end{frame}

\begin{frame}{Implications for Communication Strategy}
    \begin{cor}[Optimal Strategy Structure]
        Optimal receiver must be:
        \[
        \mathrm{Measurement} \circ \mathcal{N}
        \]
        \textbf{not}
        \[
        \mathrm{Measurement} \circ \mathcal{C} \circ \mathcal{N}
        \]
        for any CPTP $\mathcal{C}$
    \end{cor}
    
    \begin{block}{Key Insight}
        \begin{itemize}
            \item Post-channel correction cannot increase capacity
            \item Focus on optimal encoding and measurement
            \item Quantum error correction must be \textit{before} channel
        \end{itemize}
    \end{block}
\end{frame}

\section{Decoding/Inversion Ambiguity}

\begin{frame}{Channel Non-Injectivity}
    \begin{thm}[Non-Injectivity of Noisy Channels]
        For any non-unitary $\mathcal{N}$, $\exists \rho_1 \neq \rho_2$ such that:
        \[
        \mathcal{N}(\rho_1) = \mathcal{N}(\rho_2)
        \]
    \end{thm}
    
    \begin{defi}[Ambiguity Set]
        \[
        \mathcal{A}(\sigma) = \{ \rho' : \mathcal{N}(\rho') = \sigma \}
        \]
    \end{defi}
\end{frame}

\begin{frame}{Capacity-Ambiguity Trade-off}
    \begin{thm}[Capacity-Ambiguity Relation]
        \[
        C_\chi(\mathcal{N}) \leq \log_2\left( \frac{1}{\mathcal{A}_{\min}} \right)
        \]
        where $\mathcal{A}_{\min}$ is minimum achievable ambiguity
    \end{thm}
    
    \begin{block}{Connection to Week 1}
        \begin{itemize}
            \item Week 1: NP-hardness via SJIP reduction
            \item Week 2: Additional ambiguity from non-injectivity
            \item Combined: Even harder decoding problem
        \end{itemize}
    \end{block}
\end{frame}

\section{Integrated Framework: Weeks 1 \& 2}

\begin{frame}{Theoretical Progression}
    \centering
    \tiny
    \begin{tabular}{|c|p{4cm}|p{4cm}|}
        \hline
        \textbf{Stage} & \textbf{Week 1: Foundations} & \textbf{Week 2: Extensions} \\
        \hline
        \textbf{Physical} & Maxwell, quantization & GKSL, exact solutions \\
        \hline
        \textbf{Information} & Holevo capacity & Noisy capacities, DPI \\
        \hline
        \textbf{Computational} & SJIP hardness & Ambiguity characterization \\
        \hline
        \textbf{Parameters} & $\theta = \{L,\alpha,\beta_2,\ldots\}$ & $\gamma = \{\gamma_{\text{loss}},\gamma_\phi,\ldots\}$ \\
        \hline
    \end{tabular}
    
    \vspace{0.3cm}
    
    \begin{block}{Key Theorem Connections}
        \tiny
        \begin{tabular}{p{0.45\textwidth} p{0.45\textwidth}}
            \hline
            \textbf{Week 1 Theorems} & \textbf{Week 2 Extensions} \\
            \hline
            Field quantization: $\hat{A} = \sum_m \sqrt{\frac{\hbar\omega_m}{2\epsilon_0 V_m}} (\hat{a}_m f_m + \text{h.c.})$ & GKSL: $\dv{\rho}{t} = -i[H,\rho] + \sum_k \gamma_k (L_k \rho L_k^\dagger - \frac{1}{2}\{L_k^\dagger L_k, \rho\})$ \\
            \hline
            Holevo bound: $I(X:Y) \leq \chi(\{p(x), \rho(x|\theta)\})$ & Capacity non-increase: $C_\chi(\mathcal{C} \circ \mathcal{N}) \leq C_\chi(\mathcal{N})$ \\
            \hline
            SJIP NP-hardness: Subset-Sum $\leq_p$ SJIP & Non-injectivity: $\exists \rho_1 \neq \rho_2$ s.t. $\mathcal{N}(\rho_1) = \mathcal{N}(\rho_2)$ \\
            \hline
        \end{tabular}
    \end{block}
\end{frame}

\section{Progress Summary}

\begin{frame}{Week 1 Achievements}
    \begin{enumerate}
        \item Established complete electromagnetic foundation
        \item Derived quantum field quantization in fibers
        \item Developed Cq channel theory with parameter dependence
        \item Proved NP-hardness of optimal decoding via SJIP
        \item Provided explicit formulas for system design
    \end{enumerate}
\end{frame}

\begin{frame}{Week 2 Extensions}
    \begin{enumerate}
        \item Extended physical modeling to include noise via GKSL
        \item Developed exact solution methods for noisy dynamics
        \item Calculated Holevo capacities for various noisy channels
        \item Proved fundamental limits via Quantum DPI
        \item Characterized decoding ambiguity under noise
        \item Integrated all results with Week 1 foundations
    \end{enumerate}
\end{frame}

\section{Conclusion and Future Work}

\begin{frame}{Key Integrated Insights}
    \small
    \begin{enumerate}
        \item \textbf{Physical-Information Connection:}
        \[
        \gamma_{\text{loss}} = \frac{\alpha v_g \ln(10)}{10}, \quad \gamma_\phi = \frac{(\omega D_p \Delta L)^2 v_g}{2}
        \]
        
        \item \textbf{Fundamental Limits:}
        \vspace{-0.1cm}
        \begin{align*}
        I(X:Y) &\leq \chi(\{p(x), \rho(x|\theta)\}) \\
        C_\chi(\mathcal{C} \circ \mathcal{N}) &\leq C_\chi(\mathcal{N})
        \end{align*}
        
        \item \textbf{Computational Reality:} NP-hardness + ambiguity
        
        \item \textbf{Practical Guidance:} Explicit design formulas
    \end{enumerate}
\end{frame}

\begin{frame}{Future Directions}
    \begin{block}{Building on Both Weeks}
        \begin{enumerate}
            \item \textbf{Experimental validation} with fiber-optic testbeds
            \item \textbf{Algorithm development} respecting complexity constraints
            \item \textbf{System optimization} using parameter-dependent formulas
            \item \textbf{Extension} to channels with memory, nonlinear effects
        \end{enumerate}
    \end{block}
    
    \vspace{0.3cm}
    
    \begin{alertblock}{Optimization Problem}
        \[
        \max_{\theta} C(\theta) = \max_{\theta} \max_{p(x)} \chi(\{p(x), \rho(x|\theta)\})
        \]
    \end{alertblock}
\end{frame}

\begin{frame}{Final Remarks}
    \begin{block}{Complete Theoretical Picture}
        \begin{itemize}
            \item \textbf{Connects} physical reality to information limits
            \item \textbf{Bridges} ideal theory with practical considerations
            \item \textbf{Links} information limits with computational limits
            \item \textbf{Provides} explicit guidance for system design
        \end{itemize}
    \end{block}
    
    \vspace{0.5cm}
    
    \centering
    \Large
    \textbf{Integrated framework offers both theoretical \\ insights and practical engineering value}
\end{frame}

\appendix
\begin{frame}{Appendix: Mathematical Connections}
    \begin{block}{From Week 1 to Week 2 Propagation}
        Week 1 ideal:
        \[
        |\alpha(0)\rangle \to |\alpha(L)\rangle, \quad \alpha(L) = \alpha(0)e^{i\beta L}e^{-\alpha L/2}
        \]
        
        Week 2 noisy (amplitude damping):
        \[
        E_0 = \begin{pmatrix} 1 & 0 \\ 0 & \sqrt{1-\eta(t)} \end{pmatrix}, \quad E_1 = \begin{pmatrix} 0 & \sqrt{\eta(t)} \\ 0 & 0 \end{pmatrix}
        \]
        with $\eta(t) = 1 - e^{-\gamma_{\text{loss}}t}$
    \end{block}
\end{frame}

\begin{frame}{Appendix: Parameter Dependencies}
    \begin{block}{Across Weeks}
        \small
        \begin{align*}
            \text{Week 1: } & g(\theta) = d_{eg}\sqrt{\frac{\omega_0^3}{2\hbar\epsilon_0 V_{\text{eff}}(\theta)}} |f_0(r_{\text{atom}}|\theta)| \\
            \text{Week 2: } & \gamma_{\text{loss}} = \frac{\alpha v_g \ln(10)}{10}
        \end{align*}
    \end{block}
    
    \begin{block}{Connection Insight}
        Smaller $V_{\text{eff}}$ (Week 1) increases coupling but also affects loss rates (Week 2)
    \end{block}
\end{frame}

\begin{frame}
    \centering
    \Huge
    \textbf{Thank You}
    
    \vspace{1cm}
    
    \large
    Questions?
\end{frame}

\end{document}